%************************************************
\chapter{Modeling Language}\label{ch:introduction}
%************************************************
\section{A dialog between Shannon and Chomsky}
It is curious that we still don't know where the expression \textit{language model} comes from or when it became established.
We mean the \textit{expression} itself, rather than the concept.
The concept of a language model is quite old, and one of the earliest
formulations closely related to what we now call language modeling can be found in Shannon's 
work on information theory~\citep{shannon1948mathematical}\footnote{
    This intuition was already explored by Markov in 1913~\citep{1913_lm};
    Shannon later provided a more fully developed formalization.
}. 
Shannon starts from a simple observation: sequences of linguistic symbols are not random, 
but exhibit regularities characteristic of the language.
\begin{quote}
    \small\itshape
    ``[\ldots{}] the messages to be transmitted consist of sequences of letters. 
    These sequences, however, are not completely random. 
    In general, they form sentences and have the statistical structure of, say, English.''

    \begin{flushright}
    \normalfont--- Claude E. Shannon, \textit{A Mathematical Theory of Communication (1948)}
    \end{flushright}
\end{quote}
If English sentences have a statistical structure, then a statistical model that 
captures this structure should be able to generate sequences that look like English.
But what is this statistical model? How can we estimate it? Later, Shannon gives a more operational view:

\begin{quote}
    \small\itshape
    ``We can think of a discrete source as generating the message, symbol by symbol.
    It will choose successive symbols according to certain probabilities depending,
    in general, on preceding choices [\ldots{}].''

    \begin{flushright}
    \normalfont--- Claude E. Shannon, \textit{A Mathematical Theory of Communication (1948)}
    \end{flushright}
\end{quote}
Shannon thus describes a statistical
model of English but does not call it a \textit{language model}.
The story takes a different turn in the early 1970s. In 1974, a team at IBM working 
on speech recognition described a statistical model closely related to Shannon's formulation and explicitly 
called it a \textit{language model}. In their probabilistic decomposition of speech recognition, 
the language model is simply the component that assigns probabilities to possible 
word sequences---that is, a model of which utterances are more or less likely in a given language. 
As they put it, ``it is the function of the language model (LM) to provide us with $\Pr\{W \mid L\}$,'' 
where $W$ denotes a word sequence and $L$ the language. 
This work is therefore probably one of the earliest uses of the expression \textit{language model} 
in essentially its modern statistical sense: a probability distribution over linguistic sequences 
estimated from language data.

Calling these systems language models could suggest that they are
linguistic models of a language. But does a better statistical model of English mean a better model of the language itself?
The idea that statistical learning can capture linguistic structure is compatible 
with usage-based linguistics and connectionist approaches. Nativist approaches, 
however, question whether statistical learning alone is sufficient to explain linguistic knowledge.
In \textit{Three Models for the Description of Language}, Chomsky states~\citep[p.~113]{chomsky1956three}:


\begin{quote}
    \small\itshape
    ``We find that no finite-state Markov process that produces symbols with transition from state to state can serve as an English grammar. 
    [...] such processes that produce $n$-order statistical approximations to English do not come closer, with increasing $n$, to matching the output of an English grammar.''

    \begin{flushright}
    \normalfont--- Noam Chomsky, \textit{Three Models for the Description of Language (1956)}
    \end{flushright}
\end{quote}

There are two main claims made by Chomsky in this paper. 
The first claim is structural and concerns the \textbf{bounded memory} of finite-state models.
For finite-order Markov language models, this limitation manifests as a fixed context length: 
only a bounded portion of the preceding sequence can determine the next symbol.
English, however, permits recursive constructions whose dependencies can require information 
from an arbitrarily large preceding structure.
Modern language models have greatly increased the amount of preceding context they can exploit, 
making them much better at modeling long-range dependencies. 
Yet practical models still operate with bounded context and memory, 
so the distinction highlighted by Chomsky between 
unbounded linguistic structure and bounded computational memory remains relevant. 
Second, even if we consider a statistical model with unbounded memory, 
Chomsky's criticism goes beyond model capacity. His broader claim is that \textbf{language is innate},
biologically programmed in the human brain~\citep{chomsky1965aspects} and 
cannot be learned simply by accumulating statistics over observed linguistic data.
The main argument is the \textit{poverty of stimulus}: a child learns their native language rapidly,
despite being exposed to limited, noisy, and incomplete examples. According to this account, infants have structural biases that help them learn and generalize linguistic structures from a few examples.
Chomsky's theory holds that human linguistic competence depends on internal constraints 
that cannot be inferred from experience alone.

It is not straightforward, however, to translate this notion of innateness into 
the language of machine learning. Does Chomsky's argument imply that general-purpose 
learning architectures are inherently incapable of acquiring language from data alone? 
Or can we instead design computational architectures and general learning objectives whose inductive biases make
the acquisition of linguistic structure easier?
\footnote{For Chomsky, this question is not really interesting, 
    as reproducing linguistic behavior with an artificial architecture is not, 
    by itself, an explanation of human language acquisition: 
    \url{https://www.youtube.com/watch?v=PZQOIuuQiJA}}
Maybe we can also reverse the question: can learning from experience alone help
uncover the structure latent in language? Even though modern machine learning algorithms are still far from matching
the data efficiency of human children, substantial progress has been made by
designing learning architectures that optimize the language modeling objective, improving
the linguistic competence of these systems.

The language modeling objective, as Shannon describes it, is to learn to predict
the next symbol from a fixed number of previous symbols.
These symbols can be letters or words, but nowadays they are usually tokens, units that can form a word or a subword. 
More formally, the probability of picking a token $s_i$ given a sequence of previous tokens $s_{<i} = (s_1, s_2, \ldots, s_{i-1})$ is written as $P(s_i \mid s_{<i})$.
Using the chain rule, these individual probabilities can be multiplied together to give the probability of the whole text:
\begin{equation}
    P(s) = p(s_1) \times p(s_2 | s_1) \times p(s_3 | s_{1}, s_{2}) \times \cdots \times p(s_t | s_{<t}) = \prod_{i=1}^t P(s_i \mid s_{<i}).
\end{equation}
Each individual probability can be defined from a score over the vocabulary:
\begin{equation}
    p(s_i \mid s_{<i}) =
    \frac{\mathrm{score}(s_{<i}, s_i)}
    {\sum_{v \in \mathcal{V}} \mathrm{score}(s_{<i}, v)}.
\end{equation}
Here, $\mathrm{score}_\theta(x, y)$ measures how strongly the context $x$ is associated with the next token $y$, and $\mathcal{V}$ is the token vocabulary.
The function $\mathrm{score}(x, y)$ can be designed very freely but in neural language models, it is often the exponential of the dot product
between the context vector of $x$ and the token vector $y$.

Using maximum likelihood estimation, we can estimate sequence probabilities with a statistical model parameterized
by $\theta$ and trained on a large corpus $\mathcal{C}$ of texts:
\begin{equation}
    \theta^* = \arg\max_\theta
    \sum_{s \in \mathcal{C}} \sum_{i=1}^{|s|}
    \log P_\theta(s_i \mid s_{<i}).
\end{equation}

This is the learning objective of language modeling; it remains to design a learning algorithm 
that can optimize it effectively. While many neural architectures have been used for this 
purpose, ranging from CNNs to RNNs, one architecture has proven particularly effective: 
the Transformer. We next present the Transformer through the lens of linguistic 
and knowledge acquisition, and discuss how scaling this architecture leads to further gains 
in knowledge and linguistic competence.

\section{The Transformer, the best machine language learner?}
Earlier work improved neural machine translation by adding attention
to an RNN-based encoder--decoder~\citep{bahdanau2015neural}. The Transformer
replaced recurrence with attention~\citep{vaswani2017attention}. Removing recurrent dependencies allows all token positions
within a layer to be processed in parallel during training, making the architecture more
parallelizable and easier to scale. Figure~\ref{fig:decoder-only-transformer} presents the
decoder-only Transformer architecture used for autoregressive language modeling.
The model stacks $N$ layers, each containing a self-attention module followed by an MLP,
before an output head computes the next-token probability distribution. We next discuss
these three components and their roles in acquiring knowledge and generalizing linguistic
structures.

\begin{figure}[ht]
    \centering
    \includegraphics[width=0.75\linewidth]{figures/transformer.pdf}
    \caption{Decoder-only Transformer architecture.}
    \label{fig:decoder-only-transformer}
\end{figure}

\subsection{Attention}
The innovation of the Transformer is to remove the RNN and replace it entirely with
\textbf{Multi-Head Attention (MHA)}. Given $\mathbf{X} \in \mathbb{R}^{T \times d_{\mathrm{model}}}$,
the sequence of token embeddings for $T$ positions, MHA projects the sequence into queries, keys
and values for $H$ heads:
\begin{equation}
    \mathbf{Q}, \mathbf{K}, \mathbf{V}
    \in \mathbb{R}^{T \times H \times d_{\mathrm{head}}}.
\end{equation}
Equivalently, each head $h$ has its own projections:
\begin{align}
    \mathbf{Q}_h &= \mathbf{X}\mathbf{W}^{Q}_h, &
    \mathbf{K}_h &= \mathbf{X}\mathbf{W}^{K}_h, &
    \mathbf{V}_h &= \mathbf{X}\mathbf{W}^{V}_h,
\end{align}
where $\mathbf{W}^{Q}_h \in \mathbb{R}^{d_{\mathrm{model}} \times d_q}$, 
$\mathbf{W}^{K}_h \in \mathbb{R}^{d_{\mathrm{model}} \times d_k}$
and $\mathbf{W}^{V}_h \in \mathbb{R}^{d_{\mathrm{model}} \times d_v}$ are learned projection
matrices for head $h$. The head output is then computed with scaled dot-product self-attention:
\begin{equation}
    \mathrm{head}_h
    =
    \mathrm{softmax}\left(
        \frac{\mathbf{Q}_h \mathbf{K}_h^\top + \mathbf{M}}{\sqrt{d_k}}
    \right)\mathbf{V}_h .
\end{equation}
For a decoder language model, $\mathbf{M}$ is a causal mask: entries corresponding to future
positions are set to $-\infty$, so token $i$ can only attend to tokens $j \leq i$. Multi-head attention runs $H$ such heads in parallel and concatenates their outputs:
\begin{equation}
    \mathrm{MHA}(\mathbf{X})
    =
    \mathrm{Concat}\left(\mathrm{head}_1,\ldots,\mathrm{head}_H\right)\mathbf{W}^{O},
\end{equation}
with $\mathbf{W}^{O} \in \mathbb{R}^{H d_v \times d_{\mathrm{model}}}$. Each head can therefore
learn a different projection of the same context.

Having multiple independent heads is beneficial since it has been shown
that \textbf{different attention heads play linguistically interpretable roles}.
Some specialize in positional relations, syntactic
dependencies, coreference, or other linguistically interpretable patterns
\citep{clark-etal-2019-bert,voita-etal-2019-analyzing}. However, attention
heads are also substantially redundant, as many of them can be removed with
little or no degradation in model performance
\citep{michel-etal-2019-sixteen,voita-etal-2019-analyzing,
kovaleva-etal-2019-revealing}.


In standard MHA, as introduced in the original Transformer paper, the projection matrices
$\mathbf{W}^{Q}_h$, $\mathbf{W}^{K}_h$ and $\mathbf{W}^{V}_h$ are usually chosen so that all
heads together project back to a space of size $d_{\mathrm{model}}$. However, these projections
can become a bottleneck during generation. In a naive implementation, at each decoding step, all
previous tokens would be projected again to recompute their keys and values. The
\textbf{KV-Cache} reduces this computation and speeds up generation by caching the keys and values
of previous tokens. However, this cache grows linearly with the sequence length, and it is easy to
run out of memory when generating very long texts. To reduce the KV-Cache, we need to reduce the
dimensionality of keys and values.
Among the popular approaches to reduce the KV-Cache, \textbf{Grouped-Query Attention (GQA)}
shares key-value pairs across multiple query heads. Intuitively, the keys and values have fewer heads
than the queries:
\begin{align}
    \mathbf{K}, \mathbf{V} &\in \mathbb{R}^{T \times h_g \times h_{\mathrm{dim}}}, \\
    \mathbf{Q} &\in \mathbb{R}^{T \times h_q \times h_{\mathrm{dim}}},
    \qquad h_g < h_q .
\end{align}
The key-value heads are then broadcast to the query heads, and attention is computed as usual:
\begin{equation}
    \mathrm{Attention}(\mathbf{Q}, \mathbf{K}, \mathbf{V})
    =
    \mathrm{softmax}\left(\mathbf{Q}\mathbf{K}^{\top}\right)\mathbf{V}.
\end{equation}
Standard MHA corresponds to no sharing, while multi-query attention is the extreme case where all
query heads share the same keys and values. GQA is the compromise between the two: it reduces the
KV-Cache while keeping several different key-value heads.
While GQA reduces KV-cache memory, it does not reduce the attention computation. 
Nowadays, GQA is often combined with \textbf{Sparse Attention}, 
such as DeepSeek Sparse Attention~\citep{deepseekai2025deepseekv32pushingfrontieropen}, to reduce this compute 
by allowing each query to attend to only a subset of previous tokens.


\subsection{Multilayer Perceptron}
After the attention layer, the contextualized tokens are fed into the multilayer perceptron (MLP). While attention mixes information across tokens, the MLP mixes information across channels: each token is processed independently. In the original Transformer paper, the MLP is defined as:
$$
\text{MLP}(x) = f(\textbf{x}\textbf{W}_{up})\textbf{W}_{down}
$$
where $f(\cdot)$ is the non-linear activation function, 
$\textbf{W}_{\text{up}} \in \mathbb{R}^{d_{\text{model}} \times d_{\text{inter}}}$ 
projects the input token $\mathbf{x} \in \mathbb{R}^{d_{\text{model}}}$ into a higher-dimensional space,
with $d_{\text{model}} \ll d_{\text{inter}}$. Finally, $\textbf{W}_{\text{down}} \in \mathbb{R}^{d_{\text{inter}} \times d_{\text{model}}}$,
projects back to the model dimension.
Modern LLMs often replace this simple MLP with a gated linear unit variant. Instead of applying
only one non-linearity to the expanded representation, the block learns two projections: one
projection $\mathbf{V}_{\text{up}}$ produces the values, and the other projection
$\mathbf{W}_{\text{up}}$ gates them. Omitting bias terms, these alternatives can be written as:
\begin{align}
    \mathrm{MLP}_{\mathrm{GLU}}(\mathbf{x})
    &= \left(\sigma(\mathbf{x}\mathbf{W}_{\text{up}}) \odot \mathbf{x}\mathbf{V}_{\text{up}}\right)\mathbf{W}_{\text{down}}, \\
    \mathrm{MLP}_{\mathrm{GEGLU}}(\mathbf{x})
    &= \left(\mathrm{GELU}(\mathbf{x}\mathbf{W}_{\text{up}}) \odot \mathbf{x}\mathbf{V}_{\text{up}}\right)\mathbf{W}_{\text{down}}, \\
    \mathrm{MLP}_{\mathrm{Bilinear}}(\mathbf{x})
    &= \left(\mathbf{x}\mathbf{W}_{\text{up}} \odot \mathbf{x}\mathbf{V}_{\text{up}}\right)\mathbf{W}_{\text{down}}, \\
    \mathrm{MLP}_{\mathrm{SiLU}}(\mathbf{x})
    &= \left(\mathrm{SiLU}(\mathbf{x}\mathbf{W}_{\text{up}}) \odot \mathbf{x}\mathbf{V}_{\text{up}}\right)\mathbf{W}_{\text{down}}, \\
    \mathrm{MLP}_{\mathrm{ReGLU}}(\mathbf{x})
    &= \left(\mathrm{ReLU}(\mathbf{x}\mathbf{W}_{\text{up}}) \odot \mathbf{x}\mathbf{V}_{\text{up}}\right)\mathbf{W}_{\text{down}}.
\end{align}
Here, $\odot$ denotes element-wise multiplication. Table~\ref{tab:gated_mlp} gives the performance of each activation on the SuperGLUE benchmark.
The SiLU-based variant gives better results, which helps explain why many LLMs, including Mistral-v0.1 and Granite, use it.
\begin{table*}[h!]
    \centering
    \setlength{\tabcolsep}{8pt}
    \begin{tabular}{>{\raggedright\arraybackslash}p{0.30\textwidth}>{\centering\arraybackslash}p{0.40\textwidth}}
        \toprule
        \textbf{} & \textbf{Average SuperGLUE score} \\
        $\mathrm{MLP}_{\mathrm{GLU}}$ & 73.95 \\
        $\mathrm{MLP}_{\mathrm{GEGLU}}$ & 73.96 \\
        $\mathrm{MLP}_{\mathrm{Bilinear}}$ & 73.81 \\
        \rowcolor[HTML]{EEE0CC}$\mathrm{MLP}_{\mathrm{SwiGLU}}$ & \textbf{74.56} \\
        $\mathrm{MLP}_{\mathrm{ReGLU}}$ & 73.66 \\
        \bottomrule
    \end{tabular}
    \caption{Average SuperGLUE score for gated MLP variants.}
    \label{tab:gated_mlp}
\end{table*}
It is difficult to explain why gated MLPs work better with some activations than others.
As the author of the study puts it~\citep{shazeer2020gluvariants}:
\begin{quote}
    \small\itshape
    ``We offer no explanation as to why these architectures seem to work; we attribute their
    success, as all else, to divine benevolence.''

    \begin{flushright}
    \normalfont--- Noam Shazeer, \textit{GLU Variants Improve Transformer (2020)}
    \end{flushright}
\end{quote}
In modern LLM architectures, MLP parameters can account for over $90\%$ of the total model parameters.
It has been shown that MLP layers play an important role in storing knowledge acquired during web-scale pretraining. 
One interpretation treats the MLP as a key--value memory~\citep{geva-etal-2021-transformer}:
each column of \(\mathbf{W}_{\text{up}}\) in the MLP represents a key
\(\mathbf{k}_i\) that detects linguistic patterns seen during training, while
each row of \(\mathbf{W}_{\text{down}}\) represents the corresponding value
\(\mathbf{v}_i\), which promotes the tokens likely to follow those patterns.
The key \(\mathbf{k}_i\) can be strongly or weakly activated depending on the
input \(x\). When \(\mathbf{k}_i\) is strongly activated, its corresponding
value \(\mathbf{v}_i\) is assigned a large weight and therefore contributes
more strongly to the MLP output. This mechanism can be schematized as follows:
\[
\boxed{
\underbrace{\mathbf{k}_i(\mathbf{x})}_{
    \text{whether the key's pattern is active}
}
\;\longrightarrow\;
\underbrace{\mathbf{v}_i}_{
    \text{promotes tokens likely to follow the pattern}
}
}
\]
\begin{table*}[t]
    \centering
    \small
    \setlength{\tabcolsep}{6pt}
    \renewcommand{\arraystretch}{1.15}

    \begin{tabularx}{\textwidth}{@{} l l c >{\raggedright\arraybackslash}X @{}}
        \toprule
        \textbf{Value}
        & \textbf{Top prediction}
        & \textbf{Precision@50}
        & \textbf{Matching trigger example} \\
        \midrule

        $\mathbf{v}^{15}_{222}$
        & \emph{each}
        & 68\%
        & \emph{But when bees and wasps resemble \textcolor{blue}{each}} \\

        \midrule

        $\mathbf{v}^{16}_{752}$
        & \emph{played}
        & 16\%
        & \emph{Her first role was in Vijay Lalwani's psychological thriller
        Karthik Calling Karthik, where Padukone was cast as the supportive
        girlfriend of a depressed man (\textcolor{blue}{played}} \\

        \midrule

        $\mathbf{v}^{13}_{2601}$
        & \emph{extratropical}
        & 4\%
        & \emph{Most of the winter precipitation is the result of synoptic-scale,
        low-pressure weather systems (large-scale storms such as
        \textcolor{blue}{extratropical}} \\

        \midrule

        $\mathbf{v}^{15}_{881}$
        & \emph{part}
        & 92\%
        & \emph{Comet served only briefly with the fleet, owing in large
        \textcolor{blue}{part}} \\

        \midrule

        $\mathbf{v}^{16}_{2070}$
        & \emph{line}
        & 84\%
        & \emph{Sailing from Lorient in October 1805 with one ship of the
        \textcolor{blue}{line}} \\

        \midrule

        $\mathbf{v}^{12}_{3186}$
        & \emph{jail}
        & 4\%
        & \emph{On May 11, 2011, four days after scoring 6 touchdowns for the
        Slaughter, Grady was sentenced to twenty days in
        \textcolor{blue}{jail}} \\

        \bottomrule
    \end{tabularx}

    \caption{
    Examples of MLP key--value memory pairs identified in prior work
    \citep{geva-etal-2021-transformer}.
    For each value vector $\mathbf{v}^{\ell}_{i}$, \emph{Top prediction}
    is the vocabulary token receiving the highest score when the value is
    projected through the model's output embedding.
    \emph{Precision@50} is the percentage of the 50 training prefixes that
    most strongly activate the paired key $\mathbf{k}^{\ell}_{i}$ whose
    ground-truth next token equals that prediction.
    The final column presents one such matching prefix, with its
    ground-truth next token shown in blue.
    }
    \label{tab:ffn-key-value-examples}
\end{table*}
The MLP does not only store lexical knowledge. In some cases, this lexical knowledge coincides
with factual knowledge. When a key detecting the prefix \textit{The capital of France is} 
is strongly associated with a value representing \textit{Paris}, this lexical key--value memory 
turns out to also encode world knowledge. 
Hence, more broadly, it is widely believed that MLPs play a central role in storing the knowledge acquired during pretraining
\citep{geva-etal-2021-transformer,dai-etal-2022-knowledge}.
This view is also reflected in many knowledge-editing methods, which often directly target MLP weights
\citep{meng2023locatingeditingfactualassociations,meng2023masseditingmemorytransformer, fang2025alphaeditnullspaceconstrainedknowledge}.
If we view the MLP as a key--value memory, it is reasonable to expect that each input token 
requires access to only a subset of its key--value pairs. We will later discuss how this 
intuition can guide the design of an adaptive sparse memory layer.

\subsubsection{Language Modeling Head}
Well, the tokens are contextualized thanks to the multi-head attention, and each token retrieves 
the relevant knowledge encoded within the MLP. This process is repeated for $L$ layers and at the end
produces a context vector $\textbf{c}_{<i}$ that is ready for next-token prediction
using the language modeling head (LM Head or output embedding):

\begin{equation}
    p(s_i \mid s_{<i}) =
    \frac{\mathrm{\exp}(\mathbf{c}_{<i} \mathbf{W}_i)}
    {\sum_{v \in \mathcal{V}} \mathrm{\exp}(\mathbf{c}_{<i} \mathbf{W}_v)}.
\end{equation}
where $W \in \mathbb{R}^{d \times |\mathcal{V}|}$ is the LM Head weight matrix.
The LM head weight matrix is another important component of large-scale Transformer LLMs. It is closely linked to the tokenizer:
the vocabulary size determines the number of parameters in the output projection, 
while the quality of tokenization affects how easily the model learns and, ultimately, its performance.
Recent tokenizers use large vocabularies to capture the diversity of tokens in web-scale data.
Hence, the LM head can account for a large share of the parameters in an LLM, particularly in small models. Figure \ref{fig:emb-param-ratio} shows the memory
occupied by the input/output embedding parameters for different model sizes, including dense and quantized ones.
We can see that the embedding and LM Head parameters represent a larger fraction of the model 
when the Transformer backbone is small\footnote{Quantization is usually not applied to the embedding weights.}.
\begin{figure}[ht]
    \centering
    \includegraphics[width=1.0\linewidth]{figures/embedding_parameter_ratio.pdf}
    \caption{Share of model storage occupied by embedding parameters across different model sizes.
    Data taken from \href{https://blog.squeezebits.com/vocabulary-trimming-methods}{https://blog.squeezebits.com/vocabulary-trimming-methods}.}
    \label{fig:emb-param-ratio}
\end{figure}
This is mainly caused by the large vocabulary of the models. A common solution to reduce the size by half is
weight tying \citep{press-wolf-2017-using}, which shares the input and output embedding matrices. Almost all small LLMs use weight tying, which can substantially
reduce the model's parameter count without hurting performance.

Another approach to reducing the number of parameters occupied by the LM head
is to factorize its weight matrix into two low-rank matrices. This strategy was
notably adopted in ALBERT~\citep{lan2020albertlitebertselfsupervised}, where
the embedding dimension was reduced from $768$ to $128$ without sacrificing downstream performance. 
More precisely, given the input $\mathbf{x}\in\mathbb{R}^{d}$ and an LM-head weight
matrix $\mathbf{W}\in\mathbb{R}^{d\times V}$, the factorized projection is
\begin{equation}
    \mathbf{x}\mathbf{W}
    \quad\longrightarrow\quad
    (\mathbf{x}\mathbf{W}_1)\mathbf{W}_2,
\end{equation}
where $\mathbf{W}_1\in\mathbb{R}^{d\times r}$,
$\mathbf{W}_2\in\mathbb{R}^{r\times V}$, and $r\ll H$. This reduces the
LM-head parameter count from $dV$ to $dr+rV$.

However, this bottleneck architecture can hurt optimization.
Some works have identified what is known as the
\textit{softmax bottleneck}~\citep{yang2018breakingsoftmaxbottleneck}:
a low-dimensional LM head can only produce a limited variety of output
distributions. Therefore, reducing its rank too aggressively can limit the
model's ability to predict complex token distributions. A study of LM heads
with different ranks found that reducing the rank
hurts performance and limits the expressiveness of the output distribution~\citep{godey2024smalllanguagemodelsunderperform}, as shown in Figure~\ref{fig:lm_head_botlnck}.
More importantly, this also hurts the learning signal during training \citep{godey2026lostbackpropagationlmhead}.
There are many solutions to avoid the \textit{softmax bottleneck}. Since this phenomenon arises mainly when
$r \ll V$, an obvious solution is to increase $r$ or reduce the vocabulary size $V$. Reducing the vocabulary size can be
achieved using a tokenizer with fewer tokens or falling back to a character- or byte-level tokenizer.
It is also technically possible after training to remove tokens from the input and output embeddings.
But this is risky as it can affect critical tokens. This discussion is related to the research we will present
later, which proposes to reduce the rank of the LM Head instead of pruning tokens in multilingual transformer models.
\begin{figure}[ht]
    \centering
    \subfloat[Accuracy.]{
        \includegraphics[width=0.48\linewidth]{figures/figure6_accuracy.png}
    }
    \hfill
    \subfloat[Cross-entropy.]{
        \includegraphics[width=0.48\linewidth]{figures/figure6_cross_entropy.png}
    }
    \caption{Effect of reducing the LM Head rank on accuracy and cross-entropy.
    Figures from \citep{godey2024smalllanguagemodelsunderperform}.}
    \label{fig:lm_head_botlnck}
\end{figure}

\subsection{A Transformer Is Never Large Enough}
What made the Transformer architecture particularly successful in language
modeling is its ability to scale both the number of parameters and the amount
of training data. Since Transformer training can be efficiently parallelized
across multiple GPUs, increasingly large models can be trained on increasingly
large datasets, leading to improved performance and the emergence of new
capabilities. However, scaling cannot be done blindly: given the high cost of
pretraining, the allocation between model size and training data must be
estimated before the final training run. For a fixed compute budget, how large
should the model be, and on how many tokens should it be trained?

Denoting the number of parameters by $N$, the number of training tokens by $D$,
and the training compute by $C$, the compute cost of a dense Transformer can be
approximated as $C \approx 6ND$. Given a fixed training budget $C$
(the number of floating-point operations available for training), how should
we allocate compute between model size and training tokens?
Both Kaplan \citep{kaplan2020scaling} and Chinchilla \citep{hoffmann2022training} 
scaling laws indicate that the compute-optimal model size and dataset size follow power laws
in the compute budget: $N_{\mathrm{opt}} \propto C^{\alpha}$ and $D_{\mathrm{opt}} \propto C^{\beta}$.
\textbf{Kaplan scaling laws} estimate the compute-optimal allocation as
$N_{\mathrm{opt}} \propto C^{0.73}$ and  $D_{\mathrm{opt}} \propto C^{0.27}$.
Thus, model size should increase much faster than training data, 
which favors very large models trained on comparatively few tokens. 
On the other hand, the \textbf{Chinchilla scaling laws}
later found that model size and training data should instead scale approximately equally:
$N_{\mathrm{opt}} \propto C^{0.5}$ and $D_{\mathrm{opt}} \propto C^{0.5}$.

Beyond guiding compute allocation, scaling laws also reveal how model capabilities
evolve with scale. While performance often improves predictably as training compute
increases, some abilities appear only after the model reaches a certain scale. The
model may remain close to random performance across several scales before suddenly
crossing a plateau and acquiring a new ability. Several such
\textit{emergent abilities} have been identified~\citep{wei2022emergent}, as shown in Figure~\ref{fig:wei-emergent-abilities}.
\begin{figure}[ht]
    \centering
    \includegraphics[width=1.0\linewidth]{figures/wei_emergent_abilities_figure.pdf}
    \caption{Examples of emergent abilities in large language models.
    Figure from \citep{wei2022emergent}. 
    An ability is emergent when its score surpasses the random baseline
    after remaining near that baseline at smaller scales.}
    \label{fig:wei-emergent-abilities}
\end{figure}
We can see that many linguistic skills such as language understanding, especially
for low-resource languages, only show up at a given model scale.
Then a question arises. Is there a point beyond which increasing the number of parameters no longer
produces meaningful improvements? Testing this hypothesis directly is difficult,
since training models at the largest scales is extremely costly.
Classical scaling laws predict \textit{diminishing returns}: test loss continues
to decrease with scale, but each additional increase in model size produces a
smaller improvement~\citep{kaplan2020scaling,hoffmann2022training}. However,
whether this implies diminishing returns in model capabilities depends on how
performance is measured. It has been argued that short-task
benchmarks may create an illusion of saturation~\citep{sinha2026illusion}. Even small improvements in
single-step accuracy can compound into large improvements in the length of tasks
that a model can reliably execute. Therefore, diminishing returns on loss or
short-task accuracy do not necessarily imply diminishing returns in practical
capabilities.